\section{Positive kernels from repeated differences}

The logarithmically smoothed explicit formula has the arithmetic term
$\sum_n\Lambda(n)(t-\log n)_+$.  The following construction turns its
higher primitives into a nonnegative detector supported on a short
logarithmic interval.  All statements in this section are checked for
arbitrary finite order; specializing to order $128$ introduces no new
analytic argument.

\begin{definition}[Recursive primitives]
\label{kernel:primitives}
\lean{Sidon.heaviside,Sidon.heavisidePrimitive}\leanok
Put $F_0(u)=1$ for $u\geq0$ and $F_0(u)=0$ for $u<0$.
For every integer $j\geq0$, define the oriented interval integral
\[
 F_{j+1}(u)=\int_0^uF_j(v)\,dv.
\]
For a function $f$ and real step $h$, write
$\Delta_hf(u)=f(u+h)-f(u)$ and let $\Delta_h^r$ denote its $r$-fold
iteration, with $\Delta_h^0f=f$.
\end{definition}

\begin{lemma}[Integrability and the first primitive]
\label{kernel:primitive-integral}
\uses{kernel:primitives}
\lean{Sidon.intervalIntegrable_heaviside,Sidon.intervalIntegrable_heavisidePrimitive,Sidon.continuous_heavisidePrimitive_succ,Sidon.heavisidePrimitive_one,Sidon.heavisidePrimitive_succ_eq_zero_of_nonpos,Sidon.integral_heavisidePrimitive_translate,Sidon.fwdDiff_heavisidePrimitive_succ}\leanok
Every $F_j$ is integrable on every compact interval, and every $F_{j+1}$
is continuous.  Moreover $F_1(u)=\max(u,0)$, and
$F_{j+1}(u)=0$ whenever $u\leq0$.
For all real $a,b,c$,
\[
 \int_a^bF_j(v+c)\,dv
 =F_{j+1}(b+c)-F_{j+1}(a+c).
\]
In particular, for all real $h,u$,
\[
 \Delta_hF_{j+1}(u)=\int_0^hF_j(u+v)\,dv.
\]
\end{lemma}
\begin{proof}
The indicator $F_0$ of the measurable half-line $[0,\infty)$ is
measurable and bounded in absolute value by one.  A compact interval
has finite Lebesgue measure, so this constant dominates $F_0$ by an
integrable function there.  This proves interval integrability of $F_0$.
If $F_j$ is integrable on every compact interval, continuity of an
indefinite integral proves that $F_{j+1}$ is continuous.  A continuous
function is integrable on each compact interval, completing the induction.
This argument uses continuity of the integral, without asserting
differentiability of $F_0$ at zero.

If $u\geq0$, $F_0$ agrees with one on the interior of $[0,u]$, so its
integral is $u$.  If $u<0$, reverse the interval: $F_0$ is zero on
$(u,0)$ and its value at the endpoint does not change the integral.
The result is zero, proving $F_1(u)=\max(u,0)$, including $u=0$.
Inductively, if $u\leq0$ and $F_{j+1}$ vanishes on $(-\infty,0]$,
the integral defining $F_{j+2}(u)$ is zero after reversing the interval.
Thus every positive-order primitive vanishes on this half-line.

Translation invariance of Lebesgue measure changes the displayed
translated integral into $\int_{a+c}^{b+c}F_j(w)\,dw$.
The adjacent-interval identity, valid for oriented integrals and all
orders of the endpoints, rewrites this as
$\int_0^{b+c}F_j-\int_0^{a+c}F_j$.
These are exactly the two defining primitive values.
Set $a=0$, $b=h$, and $c=u$ to obtain the final formula.
No fundamental theorem of calculus is required.
\end{proof}

\begin{lemma}[Differences commute with a compact integral]
\label{kernel:difference-integral}
\lean{Sidon.intervalIntegrable_translate,Sidon.intervalIntegrable_fwdDiff_iter,Sidon.fwdDiff_iter_intervalIntegral}\leanok
Let $f:\mathbb R\to E$ take values in a normed real vector space and
be integrable on every compact interval.  For every integer $r\geq0$,
real $h$, and real $u$,
\[
 \Delta_h^r\left(x\mapsto\int_0^h f(x+v)\,dv\right)(u)
 =\int_0^h\Delta_h^rf(u+v)\,dv.
\]
All translated integrands and all finite differences appearing here are
integrable on compact intervals.
\end{lemma}
\begin{proof}
Translation preserves interval integrability, and the difference of two
integrable functions is integrable.  Induction gives this assertion for
every iterate of $\Delta_h$.
At order zero the identity is definitional.
For the induction step, evaluate the order-$r$ identity at $u+h$ and
at $u$, subtract, and use linearity of the interval integral.
Its integrand is
$\Delta_h^rf(u+h+v)-\Delta_h^rf(u+v)$.
Since $u+h+v=u+v+h$, this equals $\Delta_h^{r+1}f(u+v)$ pointwise.
The established integrability justifies the subtraction of integrals.
The proof works for real and complex integrands alike.
\end{proof}

\begin{definition}[The prime-detector kernel]
\label{kernel:detector}
\uses{kernel:primitives}
\lean{Sidon.primeDetectorKernel,Sidon.primeDetectorKernel_eq_sum}\leanok
For an integer $r\geq1$ define
\[
 K_r(h,u)=\Delta_h^rF_{r-1}(u).
\]
Equivalently, the actual finite-difference formula is
\[
 K_r(h,u)=\sum_{j=0}^r(-1)^{r-j}\binom rjF_{r-1}(u+jh).
\]
\end{definition}

\begin{lemma}[Positive compact support with an explicit bound]
\label{kernel:positivity-support}
\uses{kernel:primitive-integral,kernel:difference-integral,kernel:detector}
\lean{Sidon.primeDetectorKernel_recursion,Sidon.intervalIntegrable_primeDetectorKernel,Sidon.primeDetectorKernel_bounds,Sidon.primeDetectorKernel_eq_zero,Sidon.support_primeDetectorKernel_subset}\leanok
For $r\geq1$ and $h\geq0$, the kernel satisfies
\[
 0\leq K_r(h,u)\leq h^{r-1}\qquad(u\in\mathbb R),
 \qquad \operatorname{supp}K_r(h,\cdot)\subseteq[-rh,0].
\]
For all real $h,u$ and every integer $r\geq1$,
\[
 K_{r+1}(h,u)=\int_0^hK_r(h,u+v)\,dv.
\]
Each kernel is integrable on every compact interval.
\end{lemma}
\begin{proof}
The binomial formula follows from the finite-difference identity for
an arbitrary function.  Integrability follows because the defining
expression is a finite difference of an interval-integrable primitive.
For the recurrence, write
$\Delta_h^{r+1}F_r=\Delta_h^r(\Delta_hF_r)$.
The primitive identity expresses the inner difference as
$u\mapsto\int_0^hF_{r-1}(u+v)\,dv$.
Commuting the remaining $r$ differences through this compact integral
gives exactly the asserted recurrence.

At order one,
$K_1(h,u)=F_0(u+h)-F_0(u)$.
If $u\geq0$, both indicators equal one, so this is zero.
If $u<0$, the second indicator is zero and the first is either zero
or one.  This proves $0\leq K_1\leq1=h^0$, including $h=0$.
If $u<-h$, both indicators are zero; if $u>0$, both are one.
Thus its support is contained in $[-h,0]$.

Suppose the conclusions hold at order $r$.
The recurrence integrates a nonnegative function over $[0,h]$, so the
next kernel is nonnegative.  Its integrand is at most $h^{r-1}$;
integral monotonicity therefore bounds the result by
$h\,h^{r-1}=h^r$.
If $u>0$, then $u+v>0$ throughout $0\leq v\leq h$, so the integrand
vanishes by the induction hypothesis.
If $u<-(r+1)h$, then $u+v\leq u+h<-rh$ throughout that interval,
so the integrand again vanishes.
This proves the support assertion at order $r+1$ and completes the
induction.  Strict inequalities outside the closed support avoid any
endpoint convention issue.
\end{proof}

For the later prime application set $h=x^{-1/3}$ and
$t=\log x-rh$.  A nonzero weight $K_r(h,t-\log n)$ forces
$t\leq\log n\leq t+rh=\log x$, hence
$xe^{-rh}\leq n\leq x$ for every positive integer $n$.
Thus a positive sum of these nonnegative Mangoldt-weighted kernels
detects a prime power in an interval below $x$; removing nonprime powers
is a separate arithmetic estimate.

\begin{lemma}[Closed primitives and finite arithmetic detectors]
\label{kernel:finite-detector}
\uses{kernel:primitives,kernel:detector,kernel:positivity-support}
\lean{Sidon.heavisidePrimitive_succ_eq_pow,Sidon.fwdDiff_iter_sum_heavisidePrimitive,Sidon.sum_primeDetectorKernel_nonneg,Sidon.sum_primeDetectorKernel_le,Sidon.exists_supported_weight_of_sum_primeDetectorKernel_pos}\leanok
For $j\geq0$,
\[
 F_{j+1}(u)=\frac{\max(u,0)^{j+1}}{(j+1)!}.
\]
If $S$ is finite, $a_n,b_n$ are real, and $r\geq1$, then
\[
 \Delta_h^r\!\left(\sum_{n\in S}a_nF_{r-1}(\,\cdot-b_n)\right)(t)
 =\sum_{n\in S}a_n K_r(h,t-b_n).
\]
For $h\geq0$ and $a_n\geq0$, the latter sum lies between zero and
$h^{r-1}\sum_{n\in S}a_n$.
If it is positive, some $n\in S$ satisfies $a_n>0$ and
$t\leq b_n\leq t+rh$.
\end{lemma}
\begin{proof}
The closed form starts with $F_1(u)=\max(u,0)$.
For the induction step, first suppose $u\geq0$.
Every integration point $v\in[0,u]$ is nonnegative, so the induction
hypothesis identifies the integrand with $v^{j+1}/(j+1)!$.
The integral of this monomial is
$u^{j+2}/((j+2)(j+1)!)=u^{j+2}/(j+2)!$.
For $u\leq0$, both sides vanish by the proved support of the positive
primitives and $\max(u,0)=0$.
Forward differences are linear and commute with translation, so they
commute with the displayed finite sum and constant real multipliers.
This proves the exact identity, including order zero if desired.
Kernel positivity gives nonnegativity term by term, and its bound
$K_r\leq h^{r-1}$ gives the upper estimate after summing.
A positive finite sum of nonnegative terms has a positive term.
In that term $a_n>0$ and $K_r(h,t-b_n)>0$.
The latter forces $-rh\leq t-b_n\leq0$ by the checked kernel support,
which rearranges to the required logarithmic interval.
\end{proof}

\begin{definition}[The actual higher Mangoldt primitives]
\label{kernel:mangoldt-primitives}
\uses{kernel:finite-detector}
\lean{Sidon.smoothedMangoldtTerm,Sidon.smoothedMangoldt,Sidon.smoothedMangoldt_eq_fixed_sum,Sidon.continuous_smoothedMangoldt,Sidon.integral_smoothedMangoldt}\leanok
For integers $j\geq0$ define
\[
 A_j(t)=\sum_{1\leq n\leq\lfloor e^t\rfloor}
        \Lambda(n)F_{j+1}(t-\log n).
\]
Every $A_j$ is continuous on the whole real line, and
$\int_0^t A_j(u)\,du=A_{j+1}(t)$ for every real $t$.
In particular $A_0(t)=\sum_n\Lambda(n)(t-\log n)_+$ is the literal
arithmetic primitive appearing in the logarithmically smoothed explicit
formula.
\end{definition}
\begin{proof}
If $n=0$, its Mangoldt factor vanishes.  For $n\geq1$ outside the
moving window, $e^t<n$, hence $t<\log n$ and
$F_{j+1}(t-\log n)=0$.  Therefore for every $t\leq T$,
\[
 A_j(t)=\sum_{1\leq n\leq\lfloor e^T\rfloor}
          \Lambda(n)F_{j+1}(t-\log n).
\]
This is an exact equality: the original index set is a subset of the
larger one, and each added term is zero.  Near a fixed $t$, take
$T=t+1$.  Every summand in this fixed finite window is continuous, so
its finite sum proves continuity at $t$.

For the integral on $[0,t]$, use the common cutoff
$T=\max(t,0)$; every point of the unordered interval between $0$ and
$t$ is at most $T$.  The integrands are continuous, so integration
commutes with the finite sum.  For each $n$ the translated-primitive
identity gives
\[
 \int_0^t F_{j+1}(u-\log n)\,du
 =F_{j+2}(t-\log n)-F_{j+2}(-\log n).
\]
The second term is zero because $\log n\geq0$ for natural $n$
(including Mathlib's convention $\log0=0$).
The remaining fixed sum is $A_{j+1}(t)$ by the same cutoff equality.
This argument also treats negative $t$, using oriented interval
integrals rather than silently assuming a positive upper endpoint.
\end{proof}

\begin{lemma}[Exact positive arithmetic detector]
\label{kernel:mangoldt-detector}
\uses{kernel:mangoldt-primitives,kernel:finite-detector}
\lean{Sidon.fwdDiff_iter_smoothedMangoldt}\leanok
For $j\geq0$ and $h\geq0$,
\[
 \Delta_h^{j+2}A_j(t)=
 \sum_{1\leq n\leq\lfloor\exp(t+(j+2)h)\rfloor}
       \Lambda(n)K_{j+2}(h,t-\log n).
\]
\end{lemma}
\begin{proof}
Expand the left side by the signed binomial formula for the finite
difference.  Its arguments are $t+kh$ with $0\leq k\leq j+2$, all
bounded above by $T=t+(j+2)h$ because $h\geq0$.
Replace $A_j$ at each such argument by its exact sum in the common
window through $\lfloor e^T\rfloor$.
Finite differences commute with this finite sum, with its constant
Mangoldt coefficients, and with translation by $-\log n$.
The resulting difference of $F_{j+1}$ is $K_{j+2}$ by definition.
No infinite arithmetic series or exchange of limits is used.
\end{proof}

\begin{lemma}[Removing prime powers from the detector]
\label{kernel:prime-extraction}
\uses{kernel:mangoldt-detector,construction:prime-powers}
\lean{Sidon.exists_prime_of_kernel_mangoldt_bound,Sidon.exists_prime_of_smoothedMangoldt_bound,Sidon.nonprimeMangoldtMass_floor_isLittleO,Sidon.eventually_nonprimeKernel_error_lt}\leanok
Write
\[
 E(N)=\sum_{\substack{1\leq n\leq N\\ n\text{ not prime}}}\Lambda(n).
\]
If $r\geq1$, $h\geq0$, and
\[
 \sum_{1\leq n\leq N}\Lambda(n)K_r(h,t-\log n)
      >h^{r-1}E(N),
\]
then some actual prime $p\leq N$ satisfies
$t\leq\log p\leq t+rh$.
For $h=x^{-1/3}$, fixed $r\geq1$, and every $c>0$, eventually
\[
 h^{r-1}E(\lfloor x\rfloor)<c xh^r.
\]
\end{lemma}
\begin{proof}
Split the finite detector into its prime and nonprime summands.
Both parts are nonnegative.  The nonprime part is at most
$h^{r-1}E(N)$ because every kernel is at most $h^{r-1}$.
The strict displayed hypothesis therefore makes the prime part positive.
A positive term has prime index and nonzero kernel; kernel support then
gives the two logarithmic inequalities.

The checked prime-power estimate gives $E(N)=o(N^{2/3})$.
Compose this little-oh bound with $N=\lfloor x\rfloor$, which tends
to infinity.  Since $0\leq\lfloor x\rfloor\leq x$ eventually and
the exponent $2/3$ is positive, the denominator is at most $x^{2/3}$;
thus $E(\lfloor x\rfloor)=o(x^{2/3})$.
For any $c>0$, apply its bound with $c/2$ and use $x^{2/3}>0$
to obtain a strict bound with $c$.
Finally $x^{2/3}=x x^{-1/3}=xh$ for $x>0$.
Multiply by the positive $h^{r-1}$ and combine integer powers of $h$.
\end{proof}

\begin{definition}[Spectral smoothing]
\label{kernel:spectral-smoothing}
\uses{kernel:difference-integral,kernel:mangoldt-primitives}
\lean{Sidon.logarithmicPrimitive,Sidon.spectralSmoothing,Sidon.spectralSmoothing_recursion,Sidon.spectralSmoothing_affine,Sidon.spectralSmoothing_congr_nonneg,Sidon.logarithmicPrimitive_smoothedMangoldt,Sidon.spectralSmoothing_smoothedMangoldt}\leanok
For a continuous complex function $f$ let $P_0f=f$ and
$P_{n+1}f(t)=\int_0^tP_nf(u)\,du$, and set
$S_n(f;h,t)=\Delta_h^{n+2}P_nf(t)$.
Then
\[
 S_{n+1}(f;h,t)=\int_0^h S_n(f;h,t+u)\,du.
\]
This operation is linear, kills every affine function, and respects
identities on nonnegative time whenever $t,h\geq0$.
On the complexification of $A_0$ it is exactly the complexification of
$\Delta_h^{n+2}A_n(t)$.
\end{definition}
\begin{proof}
Successive primitives of a continuous function are continuous and hence
integrable on compact intervals.  Adjacent-integral subtraction gives
$\Delta_hP_{n+1}f(t)=\int_0^hP_nf(t+u)\,du$.
Apply the remaining $n+2$ differences and commute them with this
compact integral, obtaining the recurrence.
Linearity follows from linearity of each primitive integral and of
finite differences.  A direct two-difference calculation kills an
affine function at order zero; the recurrence then proves the claim
at every order by induction.
For an identity on nonnegative time, primitive integrals at $t\geq0$
use only points in $[0,t]$.  The binomial difference formula then uses
only $t+kh\geq0$, proving the restricted congruence statement.
Finally induction on $n$, the proved arithmetic primitive recursion,
and the fact that complexification commutes with integration identify
$P_nA_0=A_n$.  Complexification also commutes with subtraction, so it
commutes with every finite difference.
\end{proof}

\begin{lemma}[Exact spectral atom and its two bounds]
\label{kernel:spectral-atom}
\uses{kernel:spectral-smoothing}
\lean{Sidon.exponentialZeroAtom,Sidon.spectralSmoothing_exponentialZeroAtom,Sidon.spectralSmoothing_exponentialZeroAtom_norm_le,Sidon.spectralSmoothing_exponentialZeroAtom_norm_le_high,Sidon.exponential_mainTerm_lower_bound_at_log}\leanok
For $z\ne0$ put $Q_z(t)=(e^{zt}-1)/z^2$.
For $r=n+2$ one has
\[
 S_n(Q_z;h,t)=e^{zt}\left(\frac{e^{zh}-1}{z}\right)^r.
\]
If $\beta=\Re z\geq0$ and $h\geq0$, then
\[
 |S_n(Q_z;h,t)|\leq e^{\beta(t+rh)}h^r,
 \qquad
 |S_n(Q_z;h,t)|\leq
       e^{\beta(t+rh)}\frac{2^r}{|z|^r}.
\]
For the main term, $x>0$ and $rh\leq\log2$ imply
\[
 e^{\log x-rh}(e^h-1)^r\geq\tfrac{x}{2}h^r.
\]
\end{lemma}
\begin{proof}
At $n=0$, expand the two forward differences of $Q_z$; the constants
cancel and the exponential addition law factors the result into
$e^{zt}(e^{zh}-1)^2/z^2$.
For the induction step, use the spectral recurrence and the exact
translated integral
\[
 \int_0^h e^{z(t+u)}\,du=e^{zt}\frac{e^{zh}-1}{z}.
\]
The latter follows by factoring $e^{zt}$ and integrating
$e^{zu}$; division by $z$ is valid by the stated nonzero hypothesis.
The recurrence therefore multiplies the previous result by one more
factor $(e^{zh}-1)/z$.

For the low-frequency estimate, that same integral at $t=0$ gives
\[
 \left|\frac{e^{zh}-1}{z}\right|
 \leq\int_0^h e^{\beta u}\,du\leq h e^{\beta h}.
\]
For the high-frequency estimate, use
$|e^{zh}-1|\leq e^{\beta h}+1\leq2e^{\beta h}$.
Insert either bound into the exact power formula, use
$|e^{zt}|=e^{\beta t}$, and combine the real exponential factors.
Finally $e^h-1\geq h$ and $e^{rh}\leq2$ imply
$e^{\log x-rh}\geq x/2$; all factors are nonnegative, so multiplying
the two lower bounds proves the main-term inequality.
\end{proof}

\begin{lemma}[Smoothing an absolutely controlled series]
\label{kernel:smoothing-series}
\uses{kernel:spectral-smoothing}
\lean{Sidon.intervalIntegral_tsum_of_uniform_norm_bound,Sidon.logarithmicPrimitive_norm_le_of_bound,Sidon.logarithmicPrimitive_tsum_of_uniform_bound,Sidon.spectralSmoothing_tsum_of_uniform_bound}\leanok
Let $f_i:\mathbb R\to\mathbb C$ be continuous, let $m_i\geq0$ with
$\sum_i m_i<\infty$, and suppose $|f_i(u)|\leq m_i$ throughout
$0\leq u\leq T$, where $T\geq0$.
For $0\leq t\leq T$ and every $n\geq0$,
\[
 P_n\!\left(\sum_i f_i\right)(t)=\sum_i P_nf_i(t),
 \qquad |P_nf_i(t)|\leq m_iT^n.
\]
If additionally $h\geq0$ and $t+(n+2)h\leq T$, then
\[
 S_n\!\left(\sum_i f_i;h,t\right)=\sum_i S_n(f_i;h,t).
\]
All series in these statements converge absolutely at the points used.
\end{lemma}
\begin{proof}
First consider one compact interval $[0,t]\subseteq[0,T]$.
Continuity makes each $f_i$ integrable there, and
$\int_0^t|f_i(u)|\,du\leq tm_i$.
The right side is summable in $i$, so Mathlib's theorem for integration
of a series with summable absolute integrals permits exchanging the
series and this integral.

The bound on $P_n$ follows by induction.  Its base is the assumed
bound.  If $|P_nf_i(u)|\leq m_iT^n$ for $u\in[0,T]$, the norm of
its integral through $t$ is at most $tm_iT^n\leq m_iT^{n+1}$.
All factors in the last comparison are nonnegative.
This bound proves absolute convergence of $\sum_iP_nf_i(t)$ by
comparison with $T^n\sum_i m_i$.
For the identity at level $n+1$, insert the identity at level $n$
inside the primitive integral; the same absolute-integral comparison,
now with summable majorant $m_iT^n$, permits exchanging the new series
and integral.  Induction proves the first identity at every order.

Expand the final finite difference as its binomial sum of evaluations
at $t+kh$, $0\leq k\leq n+2$.
These points lie in $[0,T]$, so the proved primitive identity and
absolute convergence apply at every one.
A finite sum commutes with the absolutely convergent series, as does
multiplication by each signed binomial coefficient.
Reassemble the binomial formula term by term to obtain the spectral
smoothing identity.  This supplies an absolute-convergence proof for
every infinite-series interchange used here.
\end{proof}

\begin{lemma}[The high-zero tail needs only reciprocal-square mass]
\label{kernel:high-tail}
\uses{kernel:spectral-atom}
\lean{Sidon.weighted_spectral_atom_norm_le_high,Sidon.norm_spectral_tail_le_square_mass}\leanok
Let $z_i\in\mathbb C$, $0\leq\Re z_i\leq1$, $m_i\geq0$, and
$M=\sum_i m_i/|z_i|^2<\infty$.
If $r=n+2$, $T>0$, $h\geq0$, and $L=t+rh\geq0$, then
\[
 \left|\sum_{|z_i|\geq T}m_i S_n(Q_{z_i};h,t)\right|
 \leq e^L2^rT^{-(r-2)}M.
\]
In particular, for $L=\log x$, $h=x^{-1/3}$,
$T=x^{17/50}$ and $r=128$, its ratio to $xh^r$ is at most
$2^{128}M x^{-13/75}$, which tends to zero.
\end{lemma}
\begin{proof}
For a tail index, $|z_i|\geq T>0$ ensures $z_i\ne0$.
The high-frequency atom bound and
$e^{(\Re z_i)L}\leq e^L$ give
\[
 |m_iS_n(Q_{z_i};h,t)|
 \leq e^L2^r\frac{m_i}{|z_i|^r}
 =e^L2^r\frac1{|z_i|^{r-2}}\frac{m_i}{|z_i|^2}
 \leq e^L2^rT^{-(r-2)}\frac{m_i}{|z_i|^2}.
\]
The last inequality uses $r-2=n\geq0$ and monotonicity of nonnegative
integer powers.  Outside the tail the corresponding summand is zero,
so the same majorant applies to the series extended by zero.
Its sum is finite by the assumed mass condition.  The norm of the
series is at most the sum of its norms, and summing the displayed
comparison yields the estimate, without a numerical table of zeros.
For the displayed parameters, division by $xh^r$ leaves the exponent
\[
 \frac{128}{3}-126\frac{17}{50}=-\frac{13}{75}<0.
\]
All remaining quantities are fixed constants; hence this ratio tends
to zero.  The following nodes identify the actual zeta-zero series
and complete its density application.
\end{proof}

\begin{lemma}[Both ordinate signs and analytic multiplicity]
\label{kernel:both-signs}
\uses{construction:zero-count}
\lean{Sidon.boundedNontrivialZeros,Sidon.sum_zetaZeroOccurrences,Sidon.sum_positiveNontrivialZeros,Sidon.boundedNontrivialZeros_subset_positive_conjugates,Sidon.sum_boundedNontrivialZeros_le_twice_occurrences}\leanok
Let $Z_T$ be the finite set of actual nontrivial zeta zeros satisfying
$|\rho|\leq T$.
For every nonnegative function $f$ invariant under complex conjugation,
\[
 \sum_{\rho\in Z_T}m(\rho)f(\rho)
 \leq2\sum_{z\in\operatorname{Occurrences}(T)}f(z_1).
\]
The occurrences on the right are exactly the positive-height
multiplicity indices used in the density estimates.
\end{lemma}
\begin{proof}
The actual zero set is discrete and closed, so its intersection with
the compact closed ball is finite.  No nontrivial zero is real; the
checked eta-positivity argument proves this for the actual zeta zeros.
If $\Im\rho>0$ and $|\rho|\leq T$, then
$0<\Im\rho\leq T$, so $\rho$ belongs to the finite positive-height
set.  If $\Im\rho<0$, its literal complex conjugate belongs to that
set.  Thus $Z_T$ is contained in the union of the positive-height set
and its conjugates.
These two sets are disjoint because their imaginary parts have opposite
strict signs.  Conjugation is an involution, preserves the analytic
order of the zero, and preserves $f$; hence their weighted sums agree.
Nonnegativity permits enlarging the original finite sum to this union.
Finally expand each positive zero of analytic order $m$ into the
$m$ distinct pairs $(\rho,k)$ with $0\leq k<m$.
Pairs from different zeros are disjoint because their first coordinates
differ.  The resulting finite sum repeats $f(\rho)$ exactly $m$ times,
which proves the occurrence identity and the stated comparison.
The conjugation here is $\rho\mapsto\overline\rho$, rather than the
functional-equation symmetry $\rho\mapsto1-\overline\rho$.
\end{proof}

\begin{theorem}[The actual smoothed explicit formula]
\label{kernel:actual-smoothed-explicit}
\uses{kernel:mangoldt-detector,kernel:spectral-atom,kernel:smoothing-series}
\lean{RiemannGaussian.RosserSchoenfeldExplicitFormula.prime_eq_spectral,Sidon.smoothedMangoldt_zero_eq_primePrimitive,Sidon.smoothed_explicit_formula,Sidon.spectralSmoothing_trivialPrimitive_norm_le}\leanok
For $r=n+2$, $h,t\geq0$, the actual arithmetic detector satisfies
\[
 \begin{aligned}
 \Delta_h^r A_n(t)&=e^t(e^h-1)^r
 -\sum_{\rho}m(\rho)e^{\rho t}
       \left(\frac{e^{\rho h}-1}{\rho}\right)^r+R_n(h,t),\\
 |R_n(h,t)|&\leq\frac{\pi^2}{6}h^n.
 \end{aligned}
\]
The equality is in $\mathbb C$ after complexifying its real left
side.  The zero sum ranges over all actual nontrivial zeta zeros with
their analytic multiplicities and converges absolutely.
\end{theorem}
\begin{proof}
The imported, checked logarithmic explicit formula is
\[
 \begin{aligned}
 A_0(t)&=e^t-1-\log(2\pi)t\\
 &\quad-\sum_\rho m(\rho)\frac{e^{\rho t}-1}{\rho^2}
 +\sum_{j\geq1}\frac{1-e^{-2jt}}{4j^2}\qquad(t\geq0).
 \end{aligned}
\]
Its analytic proof, including the actual Laplace transforms, complete
xi logarithmic-derivative identity, and Fourier uniqueness argument,
is expanded in the
\href{https://github.com/TomasOrtega/Sidon/blob/master/docs/log-smoothed-explicit-formula.md}{explicit-formula proof note}.
The zero series is controlled on $[0,T]$ by the summable majorant
$(e^T+1)m(\rho)/|\rho|^2$; the trivial series is controlled there by
$1/(4j^2)$.  Apply the proved series-smoothing theorem to these two
majorants on $[0,t+rh]$.

Linearity requires care because the displayed explicit formula was
proved only at nonnegative time.  Replace each full primitive $Z(u)$
and $V(u)$ by $Z(\max(u,0))$ and $V(\max(u,0))$.
The trivial extension is continuous by continuity of $V$ on $[0,\infty)$.
For the zero extension, near any fixed $u$ choose
$T=\max(u,0)+1$.  All nearby $\max(v,0)$ lie in $[0,T]$;
the preceding summable majorant and continuity of the individual
composed atoms prove locally uniform convergence and continuity.
These globally continuous extensions agree with the original functions
at every nonnegative argument, so restricted congruence permits their
use inside $S_n$.

Apply $S_n$ to the formula.  Its affine term vanishes.
Its main term is $Q_1$, giving $e^t(e^h-1)^r$ by the exact atom identity.
Every actual zero contributes the displayed power factor, and the
proved primitive compatibility identifies the left side with
$\Delta_h^rA_n$.
For a trivial zero $z=-2j$ one has $\Re z<0$, and its term is $-Q_z$.
The integral representation gives
$|(e^{zh}-1)/z|\leq h$, while the triangle inequality gives the second
bound $|(e^{zh}-1)/z|\leq2/|z|$.
Use the first bound for $n$ factors and the second for two factors.
Since $|e^{zt}|\leq1$, this gives
$|S_n(Q_{-2j})|\leq h^n/j^2$.
The reciprocal-square sum is $\sum_{j\geq1}j^{-2}=\pi^2/6$.
Absolute summation proves the stated remainder estimate.
\end{proof}

\begin{lemma}[The detector estimate in density-count notation]
\label{kernel:detector-estimate}
\uses{kernel:actual-smoothed-explicit,kernel:both-signs,kernel:high-tail,kernel:prime-extraction}
\lean{Sidon.smoothed_zero_norm_le_occurrences_and_tail,Sidon.smoothed_detector_lower_estimate,Sidon.exists_prime_of_smoothed_total_error}\leanok
Put $L=t+rh\geq0$, $r=n+2$, $T>0$, and
$M=\sum_\rho m(\rho)/|\rho|^2$.
Then
\[
 \begin{aligned}
 e^t(e^h-1)^r&\leq\Delta_h^rA_n(t)
 +2e^Lh^r\sum_{z\in\operatorname{Occurrences}(T)}e^{-L(1-\Re z_1)}\\
 &\quad+e^L2^rT^{-(r-2)}M+\frac{\pi^2}{6}h^{r-2}.
 \end{aligned}
\]
If the three displayed errors plus
$h^{r-1}E(\lfloor e^L\rfloor)$ are strictly less than the left side,
there is an actual prime $p$ with $t\leq\log p\leq L$.
\end{lemma}
\begin{proof}
Split the actual smoothed zero series at $|\rho|<T$.
The low part has finite support in the compact zero set.
Its pointwise bound is
\[
 m(\rho)|S_n(Q_\rho)|
 \leq e^Lh^r m(\rho)e^{-L(1-\Re\rho)}.
\]
Apply the both-sign occurrence comparison to the nonnegative,
conjugation-invariant exponential weight.
The high part converges absolutely by the square-mass majorant, and
the high-tail lemma bounds its sum.
Thus the full zero series is the sum of two summable series, justifying
the split as an exact identity of ordinary absolute sums.
Take real parts in the smoothed explicit formula; the real part of a
complex number is at most its norm, and its negative real part is also
at most its norm.  Insert the two zero estimates and the trivial
remainder bound to obtain the displayed inequality.
The final strict hypothesis forces
$\Delta_h^rA_n(t)>h^{r-1}E(\lfloor e^L\rfloor)$.
The proved arithmetic prime-extraction lemma then supplies the stated
prime, including both support endpoints.
\end{proof}

\begin{theorem}[An unconditional prime interval of exponent two thirds]
\label{kernel:prime-interval}
\uses{kernel:detector-estimate,density:uniform-zeta,construction:logarithmic-reserve}
\lean{Sidon.primeDetector_high_error_eq,Sidon.primeDetector_trivial_error_eq,Sidon.eventually_prime_interval_of_zeroWeight_tendsto,Sidon.eventually_prime_interval_of_uniform_zero_density,Sidon.exists_primeIntervalBound_two_thirds}\leanok
There is a real $X$ such that for every $x\geq X$ there is a natural
prime $p$ satisfying
\[
 x-128x^{2/3}\leq p\leq x.
\]
In Lean this is the unconditional statement
$\exists X,\ \operatorname{PrimeIntervalBound}(2/3,128,X)$.
\end{theorem}
\begin{proof}
Fix the smoothing order before choosing $x$:
\[
 r=128,\quad n=r-2=126,\quad h=x^{-1/3},\quad T=x^{17/50},
\]
\[
 t=\log x-128h,\qquad s=xh^{128}>0.
\]
The variable $s$ here denotes the positive main scale, rather than a
complex zeta argument.  The checked uniform density estimate has
exponent $A=23/8$ and logarithmic loss $11$.
Since
\[
 A\frac{17}{50}=\frac{391}{400}<1,
\]
the finite-layer density estimate and the unconditional logarithmic
zero-free width give
\[
 W(x):=\sum_{z\in\operatorname{Occurrences}(x^{17/50})}
       e^{-\log x(1-\Re z_1)}\longrightarrow0.
\]
The constants in the density estimate are uniform in the real-part
threshold, so this implication applies to the moving layers near one.
The zero-free logarithmic reserve absorbs the fixed logarithmic loss.
The checked result \texttt{ZetaZeroWeight} supplies exactly this
positive-height exponential sum; no change of multiplicity convention
is necessary.

Let $M=\sum_\rho m(\rho)/|\rho|^2$, the proved finite square mass.
The high-zero and trivial-zero errors, divided by $s$, are exactly
\[
 \frac{(x2^{128}T^{-126})M}{s}
       =2^{128}M x^{-13/75},\qquad
 \frac{(\pi^2/6)h^{126}}{s}
       =\frac{\pi^2}{6}x^{-1/3}.
\]
These identities follow by multiplication of positive real powers;
the checked parameter lemmas retain the exact rational exponents.
Both right sides tend to zero.
The nonprime-power lemma gives, eventually,
$h^{127}E(\lfloor x\rfloor)<s/16$.
Also $128h\to0$, whereas $\log x\to\infty$.

Choose a single threshold $X$ so that every $x\geq X$ has $x>1$ and
all of the following properties:
\[
 W(x)<\frac1{16},\quad
 2^{128}M x^{-13/75}<\frac1{16},\quad
 \frac{\pi^2}{6}x^{-1/3}<\frac1{16},
\]
\[
 h^{127}E(\lfloor x\rfloor)<\frac{s}{16},\qquad
 128h<\min(\log2,1),\qquad \log x\geq1.
\]
The finite intersection of these eventual statements is eventual.
The final two inequalities imply $t\geq0$ and $128h\leq\log2$.
The main lower bound is therefore
$e^t(e^h-1)^{128}\geq s/2$.
The low-zero term is less than $2s/16=s/8$.
The high-zero, trivial-zero, and nonprime-power terms are each less
than $s/16$.  Their total is less than $5s/16<s/2$.
Apply the exact detector estimate and prime-extraction theorem to
obtain a prime with $t\leq\log p\leq\log x$.

Exponentiation preserves these inequalities and $p>0$, so
$e^t\leq p\leq x$.
The elementary inequality $e^{-128h}\geq1-128h$ gives
\[
 e^t=xe^{-128h}\geq x-128xh=x-128x^{2/3}.
\]
Here $xh=x^{2/3}$ uses $x>0$.
This proves the interval statement for every $x$ beyond the single
chosen threshold.  Applying the actual density existence theorem
removes all analytic premises from the final Lean declaration.
\end{proof}
